\documentclass[11pt,a4paper]{article}

% Layout
\usepackage[margin=1in]{geometry}
\usepackage{fancyhdr}
\usepackage{titlesec}
\usepackage{xcolor}
\usepackage{soul}
\usepackage{tcolorbox}

% Fonts (Computer Modern Sans Serif)
\usepackage{lmodern}
\renewcommand{\familydefault}{\sfdefault}

% Math and tables
\usepackage{amsmath, amsfonts}
\usepackage{booktabs}
\usepackage{enumitem}

% Links (no red boxes)
\usepackage[colorlinks=true,linkcolor=black,urlcolor=blue!60!black,citecolor=black]{hyperref}

% Spacing
\setlist[itemize]{noitemsep,topsep=2pt,partopsep=0pt,parsep=0pt,leftmargin=*}
\setlength{\parskip}{4pt}
\setlength{\parindent}{0pt}
\setcounter{secnumdepth}{3}

% Colors
\definecolor{accent}{HTML}{1F4E79}
\definecolor{softgreen}{HTML}{E6F4EA}
\definecolor{softyellow}{HTML}{FFF6D6}
\definecolor{softblue}{HTML}{E8F1FA}

% Section styling (larger)
\titleformat{\section}{\LARGE\bfseries\color{accent}}{\thesection}{0.8em}{}
\titleformat{\subsection}{\Large\bfseries\color{accent}}{\thesubsection}{0.8em}{}
\titleformat{\subsubsection}{\large\bfseries\color{accent}}{\thesubsubsection}{0.8em}{}

\definecolor{tocblue}{HTML}{068fc9}
\usepackage{tocloft}
\renewcommand{\cftsecfont}{\color{tocblue}}
\renewcommand{\cftsubsecfont}{\color{tocblue}}
\renewcommand{\cftsubsubsecfont}{\color{tocblue}}
\renewcommand{\cftsecpagefont}{\color{tocblue}}
\renewcommand{\cftsubsecpagefont}{\color{tocblue}}
\renewcommand{\cftsubsubsecpagefont}{\color{tocblue}}
\hypersetup{linkcolor=tocblue}

% Header / Footer
\pagestyle{fancy}
\fancyhf{}
\fancyhead[L]{\textbf{Machine Learning}}
\fancyhead[R]{\textbf{Linear Regression}}
\fancyfoot[L]{Thomaskutty Reji}
\fancyfoot[R]{\thepage}
\renewcommand{\headrulewidth}{0.4pt}
\renewcommand{\footrulewidth}{0.4pt}

% Highlight boxes
\newtcolorbox{highlightgreen}{colback=softgreen,colframe=softgreen!60!black,boxrule=0.4pt,arc=2pt,left=6pt,right=6pt,top=4pt,bottom=4pt}
\newtcolorbox{highlightyellow}{colback=softyellow,colframe=softyellow!60!black,boxrule=0.4pt,arc=2pt,left=6pt,right=6pt,top=4pt,bottom=4pt}
\newtcolorbox{highlightblue}{colback=softblue,colframe=softblue!60!black,boxrule=0.4pt,arc=2pt,left=6pt,right=6pt,top=4pt,bottom=4pt}




% Inline highlighter (no border)
\sethlcolor{yellow!20}
\newcommand{\highlighttextblue}[1]{\textcolor{blue!60!black}{#1}}
\newcommand{\highlighttextgreen}[1]{\textcolor{green!60!black}{#1}}
\newcommand{\highlighttextyellow}[1]{\hl{#1}}

% For code formatting
\usepackage{listings}
\usepackage{xcolor}

% Define Python style for listings
\lstdefinestyle{mypython}{
    language=Python,
    basicstyle=\ttfamily\small,
    backgroundcolor=\color{gray!5},
    frame=single,
    keywordstyle=\color{blue},
    commentstyle=\color{green!50!black},
    stringstyle=\color{red!60!black},
    showstringspaces=false,
    tabsize=4,
    breaklines=true
}

% Set default style for all lstlistings
\lstset{style=mypython}




\begin{document}
\begin{titlepage}
\vspace*{0.25\textheight}
{\Large\textcolor{accent}{\textbf{Data Science Concepts}}}\par
\author{Thomaskutty Reji}
\date{}

\end{titlepage}
\tableofcontents
\newpage

\section{Linear Regression}
\subsection{Simple Linear Regression}
\begin{itemize}
    \item Simple linear regression fits a straight line to data to explain how one variable (the predictor) influences another (the response).

    \item It provides both a prediction tool and an interpretation tool, showing how much the response changes for each unit change in the predictor.
\end{itemize}


\begin{align*} 
\Large Y= \beta_0 + \beta_1 X + \epsilon 
\end{align*} 


\begin{itemize}
    \item The above equation is the population regression equation;
    \item $\beta_0$ is the intercept term-that is, the expected value of $Y$ when $X=0$
    \item    
        \highlighttextgreen{$\beta_1$ is the slope—the average increase in $Y$ associated with a one unit increase in $X$.}
                
    \item The error term is a catch-all for what we miss with this simple model ( the true relationship is probably not linear, there may be other variables that cause variation in $Y$, and there may be measurement error)
    \item We typically assume that the error term is independent of $X$.

    \item The unknown coefficients $\beta_0$ and $\beta_1$ in the population regression line is unknown. We estimate these unknown coefficients using the principle of least squares.

    \begin{align*}\Large \hat{\beta_1} = \frac{\sum_{i=1}^{n} (x_i - \bar{x})(y_i-\bar{y})}{\sum_{i=1}^{n} (x_i - \bar{x})^2} = \frac{S_{xy}}{S_{xx}}
    \end{align*} 
    
    \begin{align*}
    \Large \hat{\beta_0} = \bar{y} - \hat{\beta_1}\bar{x}
    \end{align*}
    
    \item \highlighttextyellow{For each 1-unit increase in the predictor, the predicted value of the target increases by (coefficient x unit of target), holding all other variables constant.}
\end{itemize}

\subsection{Understanding Correlation and Covariance}
\begin{itemize}
    \item \textbf{Covariance} measures whether two variables move together. 
    A positive covariance means they increase or decrease together, while a negative covariance means they move in opposite directions.

    \item The \textbf{magnitude of covariance} depends on the scale of the variables, which makes it hard to compare across datasets. 
    This is why covariance alone is not a standardized measure of relationship.

    \begin{equation*}
        \Large cov(X, Y) = \frac{\sum{(X_i - \bar{X})(Y_i - \bar{Y})}}{N}
    \end{equation*}
    

    \item \textbf{Correlation} is the normalized form of covariance, obtained by dividing by the product of standard deviations. 
    This makes correlation always lie in the range $[-1, 1]$.


    \begin{align*}
        \Large Cor(X,Y) = \frac{\sum_{i=1}^{n}(x_i-\bar{x})(y_i-\bar{y})}{\sqrt{\sum(x_i-\bar{x})^2 \sum(y_i-\bar{y})^2}}
    \end{align*}

    \begin{equation*}
        \Large r = \frac{cov(X, Y)}{\sqrt{var(X) \cdot var(Y)}}
    \end{equation*}

    
    \item A correlation of \textbf{$+1$} means a perfect increasing linear relationship, 
    
    \textbf{$-1$} means a perfect decreasing linear relationship, 
    and \textbf{$0$} means no linear association.

    \item In the context of \textbf{regression}, correlation indicates the strength of the linear relationship between $X$ and $Y$.  
    Higher correlation means the regression line will fit more tightly around the data.

    \item Covariance and correlation do \emph{not} imply causation—they only describe how variables co-vary, not why.
\end{itemize}


\subsection{Multiple Linear Regression}

Multiple Linear Regression extends simple linear regression by allowing us to model a response variable $Y$ using \emph{multiple} predictors instead of just one.  
While simple regression fits a line in 2D space, multiple regression fits a hyperplane in higher dimensions, capturing more complex relationships among variables.

\[
Y = \beta_0 + \beta_1 X_1 + \beta_2 X_2 + \cdots + \beta_p X_p + \varepsilon
\]

Here, each coefficient $\beta_j$ measures the effect of predictor $X_j$ on $Y$, holding all other variables constant.

\begin{itemize}
    \item In matrix form, we write the model compactly as:
    \[
        Y = X\beta + \varepsilon,
    \]
    where $X$ is the design matrix containing all predictors (with a column of ones for the intercept).

    \item The goal is to find the parameter vector $\hat{\beta}$ that minimizes the sum of squared residuals:
    \[
        \min_{\beta} \| Y - X\beta \|^2.
    \]

    \item Solving this optimization problem gives the \textbf{normal equations}:
    \[
        X'X\hat{\beta} = X'Y.
    \]

    \item The matrix $X'X$ is known as the \textbf{Gram matrix}. It encodes all pairwise inner products between features and:
    \begin{itemize}
        \item Determines how correlated or redundant the predictors are.
        \item Must be invertible for a unique solution to exist.
        \item Can become ill-conditioned when predictors are highly collinear.
    \end{itemize}

    \item When $X'X$ is invertible, the closed-form solution for the regression coefficients is:
    \[
        \hat{\beta} = (X'X)^{-1} X'Y.
    \]

    \item This closed-form estimator is called the \textbf{Ordinary Least Squares (OLS)} solution and provides the best linear unbiased estimator (BLUE) under the Gauss–Markov assumptions.
\end{itemize}

\subsubsection{Intuition Behind $X'X$, $(X'X)^{-1}$, and $X'Y$}

\begin{itemize}
    \item \textbf{What is $X'X$ doing?}  
    \begin{itemize}
        \item $X'X$ computes all pairwise inner products between features.  
              Each entry $(X'X)_{ij}$ equals $\langle X_i, X_j \rangle$, the dot product between feature $i$ and feature $j$.  
        \item This matrix captures how similar, correlated, or redundant the predictors are.  
        \item Large off-diagonal values mean two features move together strongly (collinearity).  
        \item So $X'X$ is essentially a ``feature correlation structure'' or a geometric summary of how the predictors relate in feature space.
    \end{itemize}

    \item \textbf{What is $(X'X)^{-1}$ doing?}  
    \begin{itemize}
        \item The inverse untangles the correlations between predictors.  
        \item It removes redundancy and adjusts each coefficient so that the contribution of a predictor is measured \emph{while holding all other predictors fixed}.  
        \item If two features are highly correlated, $(X'X)^{-1}$ becomes unstable, which is why multicollinearity is dangerous: the model struggles to isolate individual effects.  
        \item Geometrically, $(X'X)^{-1}$ transforms the correlated predictor space into an orthogonal (uncorrelated) coordinate system.
    \end{itemize}

    \item \textbf{What is $X'Y$ doing?}  
    \begin{itemize}
        \item $X'Y$ computes the inner product between each feature and the response vector.  
        \item The entry $(X'Y)_j$ equals $\langle X_j, Y \rangle$, measuring how strongly predictor $j$ aligns with or explains the variation in $Y$.  
        \item This vector captures the relationship between each predictor and the target before adjusting for interactions with other predictors.
        \item In simple regression, $X'Y$ is directly related to covariance between $X$ and $Y$.
    \end{itemize}

    \item \textbf{Putting it all together:}  
    \begin{itemize}
        \item $X'Y$ tells you ``how each feature relates to $Y$''.  
        \item $X'X$ tells you ``how the features relate to each other''.  
        \item $(X'X)^{-1}$ adjusts for feature correlations so we correctly isolate each feature's unique contribution to predicting $Y$.  
    \end{itemize}
\end{itemize}

\subsection{Linear Regression Using Gradient Descent}

Gradient Descent provides an iterative optimization approach to estimate the regression coefficients instead of using the closed-form solution $(X^\top X)^{-1}X^\top y$. This is particularly useful when the number of features is very large or when matrix inversion becomes computationally expensive.

\begin{itemize}
    \item \textbf{Goal}: Minimize the Residual Sum of Squares (RSS)
    \[
        RSS = \sum_{i=1}^{n}(y_i - \hat{y}_i)^2
    \]
    where $\hat{y}_i = w^\top x_i + b$.
    
    \item \textbf{Idea}: Start with initial guesses for $w$ and $b$, then repeatedly update them in the direction that decreases the loss.  

    \item \textbf{Gradient Computation}:
    \[
        \frac{\partial RSS}{\partial w} = -\frac{2}{n} X^\top (y - \hat{y}), 
        \qquad
        \frac{\partial RSS}{\partial b} = -\frac{2}{n} \sum_{i=1}^{n}(y_i - \hat{y}_i)
    \]

    \item \textbf{Parameter Updates}:
    \[
        w_{\text{new}} = w_{\text{old}} - \alpha \frac{\partial RSS}{\partial w}, 
        \qquad
        b_{\text{new}} = b_{\text{old}} - \alpha \frac{\partial RSS}{\partial b}
    \]
    where $\alpha$ is the learning rate.

    \item \textbf{Learning Rate}: Controls the step size.  
    Too large → divergence.  
    Too small → slow convergence.

    \item \textbf{Stopping Criteria}:  
    Gradient descent continues until:
    \begin{itemize}
        \item parameters converge, or  
        \item improvement in loss becomes negligible, or  
        \item maximum iterations reached.
    \end{itemize}

    \item \textbf{Advantages}:
    \begin{itemize}
        \item Scales well to very large datasets (big $n$ and big $p$).
        \item Avoids matrix inversion.
        \item Works well for online/streaming data.
    \end{itemize}

    \item \textbf{Limitations}:
    \begin{itemize}
        \item Requires tuning of the learning rate.
        \item May converge slowly or get stuck in plateaus.
        \item Performance depends on feature scaling—standardization is often required.
    \end{itemize}

\end{itemize}

In summary, gradient descent offers a flexible and scalable way to fit linear regression, especially when closed-form solutions are computationally challenging.



\subsection{Regularization: Ridge, Lasso, and Elastic Net}

Regularization methods modify the ordinary least squares (OLS) objective by adding a penalty on the size of the coefficients.  
This helps control model complexity, reduce overfitting, and handle multicollinearity.

\begin{itemize}

    \item \textbf{Ridge Regression (L2 Regularization)}
    \[
        \hat{\beta}_{ridge} 
        = \arg\min_{\beta} 
        \left( 
            \|Y - X\beta\|_2^2 
            + \lambda \|\beta\|_2^2 
        \right)
    \]
    \begin{itemize}
        \item Penalizes the sum of squared coefficients.  
        \item Shrinks coefficients smoothly toward zero but never sets them exactly to zero.  
        \item Very useful when predictors are highly correlated.
    \end{itemize}

    \item \textbf{Lasso Regression (L1 Regularization)}
    \[
        \hat{\beta}_{lasso} 
        = \arg\min_{\beta} 
        \left( 
            \|Y - X\beta\|_2^2 
            + \lambda \|\beta\|_1 
        \right)
    \]
    \begin{itemize}
        \item Penalizes the sum of absolute values of coefficients.  
        \item Produces sparse models by forcing some coefficients exactly to zero.  
        \item Useful for feature selection.
    \end{itemize}

    \item \textbf{Elastic Net (Combination of L1 and L2)}
    \[
        \hat{\beta}_{EN} 
        = \arg\min_{\beta}
        \left(
            \|Y - X\beta\|_2^2
            + \lambda_1 \|\beta\|_1
            + \lambda_2 \|\beta\|_2^2
        \right)
    \]
    \begin{itemize}
        \item Combines Ridge and Lasso penalties.  
        \item Useful when features are correlated and we also want sparsity.  
        \item More stable than Lasso alone, especially in high-dimensional settings.
    \end{itemize}

\end{itemize}

Regularization helps prevent overfitting, improves generalization, and makes models more robust by controlling how large the coefficients are allowed to grow.
\newline
\begin{itemize}
    \item \textbf{Why reducing coefficient sizes reduces model complexity}  
    \begin{itemize}
        \item Large coefficients make the model highly sensitive to small changes in the input features, causing unstable predictions and overfitting.  
        \item By shrinking coefficients toward zero, regularization forces the model to rely only on strong, stable patterns in the data rather than noise.  
        \item Smaller weights correspond to smoother decision boundaries and simpler hypothesis functions, improving generalization.  
        \item Thus, controlling coefficient magnitude directly controls the effective complexity of the model.
    \end{itemize}
\end{itemize}

\subsection{Standard Errors of Estimates}

$$
\Large SE(\hat{\beta}_j) = \sqrt{\,RSE^2 \cdot \left( (X^\top X)^{-1} \right)_{jj} }
$$
\begin{itemize}

    \item \textbf{What the standard error represents}  
    \begin{itemize}
        \item The standard error $SE(\hat{\beta}_j)$ measures the uncertainty in the estimated coefficient $\hat{\beta}_j$.  
        \item It tells us how much $\hat{\beta}_j$ would vary if we repeatedly sampled new datasets from the same population.  
        \item A large standard error means the coefficient estimate is unstable or poorly supported by data.
    \end{itemize}

    \item \textbf{Role of the residual standard error (RSE)}  
    \begin{itemize}
        \item The term $RSE^2$ represents how much noise remains unexplained by the model.  
        \item Higher noise inflates the standard errors because it becomes harder to attribute variation in $Y$ to any specific predictor.
    \end{itemize}

    \item \textbf{Role of $(X^\top X)^{-1}$}  
    \begin{itemize}
        \item The diagonal element $\left( (X^\top X)^{-1} \right)_{jj}$ measures how much information we have about predictor $j$.  
        \item If predictors are highly correlated or do not vary much, this term becomes large, reflecting increased uncertainty.  
        \item Thus, multicollinearity directly causes inflated standard errors.
    \end{itemize}

    \item \textbf{Why we care about standard errors}  
    \begin{itemize}
        \item They allow us to compute confidence intervals for coefficients.  
        \item They are essential for hypothesis testing (e.g., is $\beta_j = 0$?).  
        \item Small standard errors mean we trust the coefficient estimates; large ones indicate that the effect may be unreliable.
    \end{itemize}

\end{itemize}

\subsection{Estimate of the sd of the error term : RSE}
The Residual Standard Error (RSE) is an estimate of the standard deviation of the error term ,$\sigma$
in the regression model:

\begin{align*}
\Large RSE = \sqrt{\frac{1}{n-p-1}RSS} = \sqrt{\frac{1}{n-p-1} \sum_{i=1}^{n} (y_i-\hat{y_i})^2}
\end{align*}


\subection{R2 Statistic}
R2 measures the proportion of variability in y that can be explained using x.


\begin{align*}
\Large R^2 = \frac{TSS-RSS}{TSS} = 1- \frac{RSS}{TSS}
\end{align*}
$$

\subsection{R\textsuperscript{2} Statistics}
The $R^2$ statistic measures the proportion of variability in the response variable $y$ that is explained by the regression model using the predictors $X$.  
It quantifies how much better the regression model is compared to simply predicting the mean of $y$.

\[
R^2 = \frac{TSS - RSS}{TSS} = 1 - \frac{RSS}{TSS}
\]

\begin{itemize}

    \item \textbf{Total Sum of Squares (TSS)}  
    \[
        TSS = \sum_{i=1}^n (y_i - \bar{y})^2
    \]
    \begin{itemize}
        \item Measures the total variability of $y$ around its mean.  
        \item Represents how spread out the data is before fitting any model.
    \end{itemize}

    \item \textbf{Residual Sum of Squares (RSS)}  
    \[
        RSS = \sum_{i=1}^n (y_i - \hat{y}_i)^2
    \]
    \begin{itemize}
        \item Measures remaining variability after fitting the model.  
        \item Small RSS means the model fits the data well.
    \end{itemize}

    \item \textbf{Explained Sum of Squares (ESS)} (Optional but useful)  
    \[
        ESS = \sum_{i=1}^n (\hat{y}_i - \bar{y})^2
    \]
    \begin{itemize}
        \item Measures how much of the total variability the model is able to explain.
    \end{itemize}

    \item \textbf{Interpretation of $R^2$}
    \begin{itemize}
        \item $R^2 = 0$ means the model explains none of the variability—no better than predicting the mean.  
        \item $R^2 = 1$ means perfect prediction.  
        \item Higher $R^2$ means the fitted regression line captures more of the variation in $y$.
    \end{itemize}

    \item \textbf{Graphical intuition (see plot above)}
    \begin{itemize}
        \item The horizontal dashed line shows $\bar{y}$; deviations from it form the TSS.  
        \item The regression line explains part of this variability (ESS).  
        \item The leftover vertical deviations from the regression line are the RSS.  
        \item $R^2$ simply tells us:  
        \[
            \text{How much of the total deviation from the mean is eliminated by using the regression model?}
        \]
    \end{itemize}

\end{itemize}


\subsection{Adjusted R\textsuperscript{2}}
Adjusted $R^2$ is a modification of the usual $R^2$ that accounts for the number of predictors used in the model.  
While $R^2$ always increases (or at least never decreases) when more variables are added,  
Adjusted $R^2$ increases only if a new predictor improves the model more than would be expected by chance.

\begin{align*}
    \Large \text{Adjusted} R^2 
    =  1 - \frac{(1 - R^2)(N - 1)}{N - p - 1}
\end{align*}


\begin{itemize}

    \item \textbf{Why we need Adjusted $R^2$}
    \begin{itemize}
        \item Adding a new predictor always reduces RSS, which makes $R^2$ artificially increase—even if the predictor is useless.
        \item Adjusted $R^2$ corrects this by introducing a penalty for adding more features.
        \item It allows fair comparison between models with different numbers of predictors.
    \end{itemize}

    \item \textbf{How the penalty works}
    \begin{itemize}
        \item The denominator $N - p - 1$ shrinks as the number of predictors $p$ grows.
        \item This makes the fraction larger, which decreases Adjusted $R^2$ when unnecessary predictors are added.
        \item Only predictors that substantially reduce RSS will cause Adjusted $R^2$ to increase.
    \end{itemize}

    \item \textbf{Interpretation}
    \begin{itemize}
        \item Higher Adjusted $R^2$ indicates a model that explains variability efficiently, without relying on too many predictors.
        \item If Adjusted $R^2$ decreases when a new variable is added, that variable is likely not contributing useful information.
    \end{itemize}

\end{itemize}


\subsection{AIC & BIC Calculation}

\begin{itemize}
    \item Both AIC and BIC are model selection criteria designed to choose the best model among several candidates.
    \item They reward models with good fit (low RSS) but penalize models with too many parameters.
    \item The key idea: \textbf{a good model explains the data well using as few parameters as possible}.
    \item AIC is more focused on prediction accuracy, often choosing slightly more complex models.
    \item BIC imposes a stronger penalty for complexity, often leading to simpler models, especially as sample size $n$ grows.
    \item Both are especially useful when comparing non-nested models (models that are not subsets of each other).
    \item Lower AIC/BIC values indicate a better model.
\end{itemize}

\[
AIC = n \log\left(\frac{RSS}{n}\right) + 2k
\]

\[
BIC = n \log\left(\frac{RSS}{n}\right) + k\log(n)
\]


\subsection{Hypothesis Testing of Feature Significance}

In linear regression, we often want to test whether a particular predictor $X_1$ has a meaningful effect on the response variable $Y$.  
This is done through a hypothesis test on its coefficient $\beta_1$.

\[
H_0 : \beta_1 = 0 \quad \text{(no relationship between $X_1$ and $Y$)}
\]

\[
H_1 : \beta_1 \ne 0 \quad \text{(some relationship exists)}
\]

To test this, we compute the $t$–statistic:

\[
\Large 
t = \frac{\hat{\beta}_1 - 0}{SE(\hat{\beta}_1)}
\]

\begin{itemize}
    \item The numerator measures how far the estimated coefficient is from the null hypothesis value.
    \item The denominator (standard error) reflects the uncertainty in the estimate.  
          Large uncertainty makes it harder to claim significance.
    \item A large $|t|$ indicates that $\hat{\beta}_1$ is far from zero relative to its noise level.
    \item We compare this $t$ value with a $t$-distribution (with $n - p - 1$ degrees of freedom) to compute a p-value.
    \item A small p-value means the observed effect is unlikely under $H_0$, leading us to reject the null hypothesis.
\end{itemize}

This test helps determine whether a predictor contributes meaningful explanatory power to the model.



\subsection{F - test in regression analysis}
\begin{itemize}
\item Is at least one of the predictors $X_1,X_2,...,X_p$ useful in predicting the response?
\item Do all the predictors help to explain $Y$, or is only a subset of the predictors useful?
\item $H_o : \beta_1 = \beta_2 = ... = \beta_p = 0$ versus the alternative
\item $H_a :$ at least one $\beta_j$ is non-zero.

\begin{align*}
\Large F = \frac{(TSS-RSS)/p}{RSS/(n-p-1}
\end{align*}

\item IF f statistic takes a value close to 1, then we cannot expect any relationship between the response and predictors.
\item Larger the f value higher the the evidence against the null hypothesis.

\item Sometimes, we want to test that a particular subset of q of the coefficients are zero.
\item In this case we fit two model and compute their respective RSS value(with and without that particular subset).

\begin{align*}
\Large F = \frac{ (RSS_{\text{reduced}} - RSS_{\text{full}}) / q }{ RSS_{\text{full}} / (n - p - 1) }
\end{align*}

\end{itemize}


\subsection{Confidence Intervals \& Prediction Intervals}

Confidence intervals and prediction intervals are essential tools in regression analysis for quantifying the uncertainty in estimates and predictions.  
While point estimates (e.g., $\hat{\beta}_j$ or $\hat{y}$) give single best guesses, these intervals express how confident we are about those estimates based on the available data.

\subsubsection{Confidence Intervals for Regression Coefficients}

For each coefficient $\beta_j$ in the linear model
\[
Y = X\beta + \varepsilon, \quad \varepsilon \sim N(0, \sigma^2 I),
\]
we can construct a confidence interval (CI) for the estimated coefficient $\hat{\beta}_j$.

\[
\hat{\beta}_j \pm t_{\alpha/2,\, n-p-1} \cdot SE(\hat{\beta}_j)
\]

\begin{itemize}
    \item $SE(\hat{\beta}_j)$ is the standard error of the estimate.
    \item $t_{\alpha/2,\, n-p-1}$ is the critical value from the $t$-distribution with $n - p - 1$ degrees of freedom.
    \item A 95\% CI means: \\
    If we repeatedly sampled datasets and refit the model, 95\% of such intervals would contain the true $\beta_j$.
\end{itemize}

\textbf{Interpretation:}
\begin{itemize}
    \item If the CI contains zero, the predictor may not be statistically significant.
    \item Narrow CIs indicate precise estimates; wide CIs indicate high uncertainty.
\end{itemize}

\subsubsection{Confidence Interval for the Mean Response}

For a new input $x_0$ (a row vector of predictors), the expected response is:
\[
\hat{y}_0 = x_0^\top \hat{\beta}
\]

A $(1-\alpha)$ confidence interval for the \emph{mean response} at $x_0$ is:

\[
\hat{y}_0 \pm t_{\alpha/2,\, n-p-1} \cdot 
RSE \sqrt{x_0^\top (X^\top X)^{-1} x_0}
\]

\begin{itemize}
    \item This interval quantifies the uncertainty in estimating the \emph{average} value of $Y$ at $x_0$.
    \item It does not account for individual outcome noise—only uncertainty in estimating the regression line.
\end{itemize}

\subsubsection{Prediction Interval for a New Observation}

A prediction interval (PI) accounts for two sources of uncertainty:

\begin{enumerate}
    \item Uncertainty in estimating the regression line.
    \item The inherent noise in individual observations ($\varepsilon$).
\end{enumerate}

The formula for a new predicted value $Y_0$ at $x_0$ is:

\[
\hat{y}_0 \pm t_{\alpha/2,\, n-p-1} \cdot 
RSE \sqrt{1 + x_0^\top (X^\top X)^{-1} x_0}
\]

\textbf{Key differences from confidence intervals:}

\begin{itemize}
    \item The extra ``1'' under the square root accounts for the random noise in future observations.
    \item Prediction intervals are always wider than confidence intervals.
\end{itemize}

\subsubsection{Why These Intervals Matter}

\begin{itemize}
    \item They quantify the uncertainty of estimates and predictions.
    \item They help assess the reliability of model-based decisions.
    \item They allow hypothesis testing and significance evaluation.
    \item They are essential in forecasting scenarios where predicting future outcomes requires uncertainty quantification.
\end{itemize}

\subsubsection{Summary}

\begin{itemize}
    \item \textbf{Coefficient CI:} Uncertainty in $\hat{\beta}_j$.  
    \item \textbf{Mean Response CI:} Uncertainty in the average value $\mathbb{E}[Y|x_0]$.  
    \item \textbf{Prediction Interval:} Uncertainty in an individual new observation $Y_0$ at $x_0$.  
    \item \textbf{Width relationship:} \\
    Prediction Interval $\;>\;$ Confidence Interval for Mean $\;>\;$ Point Estimate.
\end{itemize}

These tools collectively allow us to understand not just the fitted model,  
but also the reliability of its estimates and predictions.

\subsection{OLS Violations}

Ordinary Least Squares relies on several key assumptions.  
Violations of these assumptions can lead to unstable coefficients, incorrect inference,  
or inefficient predictions. Below are the main types of OLS assumption violations.

\subsubsection{Multicollinearity}

\begin{itemize}
    \item Occurs when two or more explanatory variables are highly correlated.
    \item Multicollinearity increases the variance of coefficient estimates, making confidence intervals wider.
    \item The model becomes unstable: small changes in the data can cause large changes in coefficients.
    \item Diagnosis: Variance Inflation Factor (VIF) is the conventional tool.
    \item Warning signs:
    \begin{itemize}
        \item High overall $R^2$ but individual predictors have high p-values (insignificant).
        \item Coefficient signs differ from theoretical expectations.
        \item Adding or removing a feature causes coefficients to change dramatically.
        \item Adding or deleting observations changes coefficients substantially.
    \end{itemize}
    \item Remedies:
    \begin{itemize}
        \item Combine correlated variables (feature engineering).
        \item Remove or replace redundant features.
        \item Use regularization (Ridge, Elastic Net).
    \end{itemize}
\end{itemize}

\subsubsection{Heteroskedasticity}

\begin{itemize}
    \item Refers to non-constant variance of the error term across observations.
    \item Residuals exhibit unequal scatter across the range of fitted values.
    \item Causes inefficient and biased standard errors, leading to unreliable hypothesis tests.
    \item Detection:
    \begin{itemize}
        \item Residuals vs. fitted plot (visual pattern or “funnel shape”).
        \item Breusch–Pagan test (formal test).
    \end{itemize}
    \item Remedies:
    \begin{itemize}
        \item Transformations (log, square root).
        \item Robust (heteroskedasticity-consistent) standard errors.
        \item Weighted least squares.
    \end{itemize}
\end{itemize}

\subsubsection{Non-Normality of Errors}

\begin{itemize}
    \item Violates the assumption that residuals are normally distributed.
    \item Affects confidence intervals, p-values, and significance tests.
    \item Detection:
    \begin{itemize}
        \item Q-Q plot.
        \item Shapiro–Wilk, Kolmogorov–Smirnov, Anderson–Darling tests.
    \end{itemize}
    \item Common sources:
    \begin{itemize}
        \item Outliers.
        \item Skewed or heavy-tailed distributions.
        \item Omitted variables or incorrect functional form.
    \end{itemize}
    \item Remedies:
    \begin{itemize}
        \item Box–Cox or log transformations.
        \item Remove or address outliers.
        \item Respecify the model (modify functional form).
    \end{itemize}
\end{itemize}

\subsubsection{Serial Correlation (Autocorrelation)}

\begin{itemize}
    \item Occurs when error terms are correlated across time (common in time series data).
    \item Violates independence of errors, leading to underestimated standard errors.
    \item Detection:
    \begin{itemize}
        \item Durbin–Watson test.
        \item Plot of residuals over time.
    \end{itemize}
    \item Remedies:
    \begin{itemize}
        \item Include lagged variables.
        \item Use time-series models (ARIMA, GLS).
        \item Apply Newey–West robust standard errors.
    \end{itemize}
\end{itemize}

\subsection{Residual Diagnostics}

Residual diagnostics help verify whether the OLS assumptions hold for a fitted regression model.  
These tools allow us to identify non-linearity, heteroskedasticity, non-normal errors, serial correlation,  
influential observations, and multicollinearity.

%-----------------------------------------------------------
\subsubsection{Residuals vs Fitted Plot}

\begin{itemize}
    \item Used to assess non-linearity and heteroskedasticity.
    \item Ideally, residuals should be randomly scattered around zero with no visible pattern.
    \item Systematic curves suggest non-linearity; changing spread suggests heteroskedasticity.
\end{itemize}

%-----------------------------------------------------------
\subsubsection{Q--Q Plot (Quantile--Quantile)}

A Q--Q plot compares the quantiles of the residuals to those of a theoretical normal distribution.

\begin{itemize}
    \item Points on a straight line indicate normal residuals.
    \item Curvature suggests skewness.
    \item Heavy tails appear as deviations at the ends of the plot.
    \item Outliers appear far from the reference line.
\end{itemize}

%-----------------------------------------------------------
\subsubsection{Shapiro--Wilk Test for Normality}

\begin{itemize}
    \item Tests whether residuals follow a normal distribution.
    \item Null hypothesis: residuals are normally distributed.
    \item Particularly powerful for small sample sizes ($n < 50$).
\end{itemize}

%-----------------------------------------------------------
\subsubsection{Kolmogorov--Smirnov Test for Normality}

\begin{itemize}
    \item Compares the empirical distribution of residuals to a theoretical normal CDF.
    \item Null hypothesis: residuals follow the normal distribution.
    \item Sensitive to differences in both shape and location.
\end{itemize}

%-----------------------------------------------------------
\subsubsection{Anderson--Darling Test for Normality}

\begin{itemize}
    \item A refinement of the K--S test with extra sensitivity to tail behavior.
    \item Null hypothesis: residuals are normally distributed.
    \item Strong at detecting heavy tails or deviations in extremes.
\end{itemize}

%-----------------------------------------------------------
\subsubsection{One-Sample t-test: Mean of Residuals = 0}

\begin{itemize}
    \item Tests whether the residuals have zero mean, as required under OLS assumptions.
    \item Null hypothesis: $\mathbb{E}[e_i] = 0$.
\end{itemize}

\[
t = \frac{\bar{e}}{s_e / \sqrt{n}},
\qquad
\bar{e} = \text{mean of residuals},\quad
s_e = \text{standard deviation of residuals},\quad
n = \text{sample size}.
\]

%-----------------------------------------------------------
\subsubsection{Leverage and Influence (Leverage Plot)}

\begin{itemize}
    \item Leverage measures how far an observation's predictor values are from the mean of the predictors.
    \item High-leverage points can strongly affect the regression surface.
    \item Even a point with a small residual can be influential if its leverage is high.
\end{itemize}

\textbf{Cook's Distance}:
\begin{itemize}
    \item Measures the overall influence of an observation on the fitted regression coefficients.
    \item Rule of thumb: $D_i > 1$ indicates potentially influential observations.
\end{itemize}

%-----------------------------------------------------------
\subsubsection{Scale--Location Plot}

\begin{itemize}
    \item Also known as the Spread--Location plot.
    \item Plots $\sqrt{|e_i|}$ against fitted values to diagnose heteroskedasticity.
    \item A horizontal line with random scatter suggests constant variance.
    \item Systematic patterns indicate heteroskedasticity:
    \begin{itemize}
        \item Upward trend: variance increases with fitted values (classic funnel shape).
        \item Downward trend: variance decreases with fitted values.
    \end{itemize}
\end{itemize}

%-----------------------------------------------------------
\subsubsection{Breusch--Pagan Test for Homoscedasticity}

\begin{itemize}
    \item Tests whether error variance is constant.
    \item Null hypothesis: homoscedastic errors.
    \item Alternative: heteroskedasticity dependent on predictors.
\end{itemize}

\[
\hat{\varepsilon}_i^2
\]

The auxiliary regression is:

\[
\hat{\varepsilon}_i^{\,2} = \gamma_0 + \gamma_1 x_{i1} + \cdots + \gamma_p x_{ip} + u_i
\]

Let $R_{\text{aux}}^{2}$ be the $R^2$ from this regression.  
The Breusch--Pagan statistic is:

\[
W = n R_{\text{aux}}^{2},
\qquad W \sim \chi^2_{p} \text{ under the null}.
\]

%-----------------------------------------------------------
\subsubsection{Durbin--Watson Test for Autocorrelation}

\begin{itemize}
    \item Detects first-order autocorrelation in residuals (commonly in time-series data).
    \item Null hypothesis: no first-order autocorrelation.
\end{itemize}

Correct formula:

\[
d = \frac{\sum_{t=2}^{T} (e_t - e_{t-1})^{2}}{\sum_{t=1}^{T} e_t^{2}}.
\]

\begin{itemize}
    \item $d \approx 2$: no autocorrelation.
    \item $d < 1.5$: positive autocorrelation.
    \item $d > 2.5$: negative autocorrelation.
\end{itemize}

%-----------------------------------------------------------
\subsubsection{Variance Inflation Factor (VIF): Multicollinearity}

\[
VIF_j = \frac{1}{1 - R_j^{2}},
\]

where $R_j^{2}$ is obtained from regressing predictor $X_j$ on all other predictors.

\begin{itemize}
    \item $VIF = 1$: no multicollinearity.
    \item $VIF > 5$: moderate multicollinearity concern.
    \item $VIF > 10$: severe multicollinearity.
\end{itemize}

\newpage
\section{Logistic Regression}





\end{document}
